\section{Background}
\label{sec:background}

\subsection{Audit Logging in Healthcare Systems}

ISO/IEC 27789 defines the hierarchy of audit artifacts: an \emph{audit
event} triggers one \emph{audit record}; a chronological collection of
records forms an \emph{audit log}; one or more logs collectively constitute
an \emph{audit trail}~\cite{iso27789_2021}. A well-formed audit log must
satisfy four properties: (i)~\emph{unambiguous identification}---sufficient
data to identify every user or external system that accessed health data;
(ii)~\emph{integrity and confidentiality}---protection against tampering,
including logging events when the audit system itself is disabled; (iii)~%
\emph{completeness}---all create, read, update, and archive operations on
health data must be recorded; and (iv)~\emph{availability}---entries must
be captured whenever the health information system is
operational~\cite{iso27789_2021}.

NIST SP 800-53 Rev.\ 5 further formalizes these requirements through the
Audit and Accountability (AU) control family, covering event definition
(AU-2), record content (AU-3), log protection (AU-9), and retention
(AU-11)~\cite{nist80053r5}. Audit logs are also the primary evidentiary
source for detecting and attributing insider threats---a particularly
relevant concern in multi-role healthcare environments where actors hold
legitimate access by default~\cite{cert2022insiderthreat}.

\subsection{The DecMed Framework}

DecMed is a three-tier decentralized EHR framework~\cite{DecMed}. Its
\emph{client layer} comprises three Tauri/SvelteKit applications: Patient
Client (mobile), Hospital Client (desktop), and Ministry Client (desktop).
All cryptographic identity is managed non-custodially via BIP39 mnemonic
phrases stored on user devices.

The \emph{proxy layer} centers on a Rust/Axum PRE Server that mediates
data access. When a healthcare provider requests a patient record, the PRE
Server (1) verifies an active capability entry on the IOTA smart contract,
(2) re-encrypts the IPFS-stored clinical ciphertext under the requester's
public key using pre-uploaded key fragments (\texttt{kfrag}), and (3)
returns the re-encrypted payload. A Gas Station Server sponsors IOTA
transaction fees so end-users need not hold IOTA tokens.

The \emph{storage layer} combines IOTA smart contracts (Move language)
for on-chain state---capability entries, administrative metadata---and IPFS
Cluster for off-chain encrypted clinical data. IOTA Rebased uses a DAG
consensus (Mysticeti, BFT Delegated Proof-of-Stake) achieving approximately
500\,ms finality at over 50{,}000 TPS.

DecMed claims to produce tamper-evident audit trails by recording
capability entries on-chain~\cite{DecMed}. However, capability entries
represent \emph{granted permissions}, not realized accesses; revocation
overwrites state rather than appending; and event coverage is limited to
smart-contract invocations, missing interactions that do not trigger a Move
call.

\subsection{STRIDE Threat Modeling}

STRIDE is a threat classification framework developed at
Microsoft~\cite{shostack2014}. It enumerates six categories:
\textbf{S}poofing (false identity), \textbf{T}ampering (unauthorized data
modification), \textbf{R}epudiation (deniability of an action),
\textbf{I}nformation Disclosure (unauthorized exposure), \textbf{D}enial of
Service (availability degradation), and \textbf{E}levation of Privilege
(exceeding authorization). Applied per-interaction on a Data Flow
Diagram~\cite{owasp_threat_modeling}, each data flow crossing a trust
boundary is evaluated against applicable categories. The Repudiation
category has a direct correspondence with audit logging: if a system lacks,
fails to retain, or cannot verify its logs, repudiation threats cannot be
rebutted~\cite{shostack2014}.

\subsection{Hash-Chaining for Tamper-Evidence}

Hash-chaining is a lightweight integrity mechanism in which each log record
includes the SHA-256 digest of the immediately preceding record. A chain
of $N$ records is tamper-evident: modifying, inserting, or deleting any
record invalidates all subsequent hashes, making manipulation detectable
without requiring trusted hardware~\cite{nist800_92}.

\subsection{IPFS and Content-Addressed Storage}

IPFS is a peer-to-peer distributed file system that addresses data by a
cryptographic hash of its content, called a Content Identifier
(CID)~\cite{benet2014ipfscontentaddressed}. Because a CID uniquely
identifies a specific byte sequence, any modification produces a distinct
CID. IPFS nodes apply \emph{garbage collection} (GC) to remove
unpinned blocks; pinning marks a CID as permanent, preventing GC from
removing it.